\chapter{Schwierigkeiten}

%Sollte vorerst passen, wird aber noch ausgebaut, mit möglichen Lösungen/unseren Lösungen, weitere Schwieriegkeiten was die Daten angeht
Bei der späteren Erfassung könnte es später zu Schwierigkeiten kommen, diese sind mitunter die verschiedene Konditionen des menschlichen Körpers sowie der Umgebung als auch das anbringen der Hardware an dem menschlichen Körper.

\section{Einfluss von Testpersonen}
Obwohl das Gehen eine typisch menschliche Fortbewegungsart ist, variiert jeder Bewegungsablauf individuell stark. Diese Unterschiede werden von persönlichen Faktoren wie Konstitution, Beweglichkeit, Kraft, Koordinationsfähigkeit der beteiligten Muskeln sowie Stimmung beeinflusst \cite{suppe2013fbl}. Hinzu kommen Körperbau, Beweglichkeit und Gangbild. Es stellt sich also die Schwierigkeit, ein einheitliches System zu entwickeln, das auf mehrere Personen anwendbar ist. Ein solches System muss in der Lage sein, die individuellen Unterschiede zu berücksichtigen und dennoch vergleichbare und reproduzierbare Ergebnisse zur Erkennung der Bewegung zu liefern. In Bezug auf dieses Thema stellen sich zwei Fragen: Wie können Verzerrungen vermieden werden und welche Variation ist notwendig?

\subsection{Vermeidung von Verzerrungen}
Eine der zentralen Herausforderungen bei der Entwicklung eines Bewegungserkennungssystems ist die Vermeidung von Verzerrungen, die durch unbeabsichtigte Bewegungen oder äußere Einflüsse entstehen können. Um die Genauigkeit und Zuverlässigkeit der gesammelten Daten zu gewährleisten, müssen standardisierte Testbedingungen geschaffen werden. Ein wichtiger Aspekt dabei ist die Verwendung einheitlicher Testprotokolle. Diese Protokolle sollten detaillierte Anweisungen enthalten, wie die Testpersonen sich während der Tests verhalten sollen, um konsistente und vergleichbare Ergebnisse zu erzielen. Beispielsweise sollten die Testpersonen darauf hingewiesen werden, die Tasche mit dem Arduino nicht mit den Armen zu berühren, um ungewollte Bewegungen zu vermeiden, die die Daten verfälschen könnten. Zusätzlich sollten kontrollierte Umgebungsbedingungen geschaffen werden. Dies kann durch die Verwendung eines standardisierten Testbereichs erreicht werden, in dem äußere Einflüsse wie unebene Oberflächen oder Hindernisse minimiert werden. Ein Beispiel hierfür wäre die Verwendung eines Laufbands mit festgelegter Geschwindigkeit und Neigung, um sicherzustellen, dass alle Testpersonen unter denselben Bedingungen getestet werden.

\subsection{Notwendige Variation}
Ein weiterer wichtiger Aspekt ist die Frage, wie viel Variation erforderlich ist, um aussagekräftige Ergebnisse zu erzielen. Eine zu geringe Variation könnte dazu führen, dass wichtige Unterschiede zwischen den Testpersonen nicht erfasst werden. Eine zu große Variation könnte hingegen die Vergleichbarkeit der Ergebnisse beeinträchtigen. Um genügend Variation einzubringen, wurde entschieden, einen gleichen Anteil an Frauen und Männern unterschiedlicher Größe als Testsubjekte zu wählen. Dazu wird eine Pilotstudie durchgeführt, um die optimale Variation zu bestimmen. Auf Basis dieser Daten können dann die endgültigen Testbedingungen festgelegt werden, die eine ausreichende Variation gewährleisten, ohne die Vergleichbarkeit der Ergebnisse zu gefährden.

\section{Umwelteinflüsse}
Bei der Erfassung und Analyse von Bewegungsdaten spielen Umwelteinflüsse eine entscheidende Rolle. Verschiedene Umweltfaktoren können die Genauigkeit und Zuverlässigkeit der gesammelten Daten beeinflussen. Im Folgenden werden einige der wichtigsten Umwelteinflüsse erläutert und abschließend eine Lösung präsentiert.

\begin{itemize}
	\item \textbf{Bodenbeschaffenheit}: Die Beschaffenheit des Bodens, auf dem sich die Testperson bewegt, kann erhebliche Auswirkungen auf die Bewegungsmuster haben. Unterschiedliche Oberflächen wie Asphalt, Gras, Sand oder unebener Untergrund können die Art und Weise, wie eine Person geht oder läuft, beeinflussen. Beispielsweise erfordert das Gehen auf Sand mehr Muskelkraft und führt zu anderen Bewegungsmustern als das Gehen auf einer festen, ebenen Oberfläche.
	\item \textbf{Umgebung}: Die Umgebung, in der die Tests durchgeführt werden, kann ebenfalls einen Einfluss haben. Enge Räume, Hindernisse oder belebte Umgebungen können die Bewegungsfreiheit einschränken und zu unnatürlichen Bewegungsmustern führen. Es ist daher wichtig, eine kontrollierte und standardisierte Testumgebung zu schaffen, um vergleichbare Ergebnisse zu erzielen.
	\item \textbf{Bewegungsziel}: Das Ziel der Bewegung kann ebenfalls eine Rolle spielen. Beispielsweise kann das Gehen zu einem bestimmten Ziel schneller und direkter sein als das Gehen ohne spezifisches Ziel. Dies kann zu unterschiedlichen Bewegungsmustern führen, die bei der Analyse berücksichtigt werden müssen.
	\item \textbf{Temperatur}: Die Umgebungstemperatur kann die Muskelaktivität und die allgemeine Leistungsfähigkeit der Testperson beeinflussen. Extreme Temperaturen, sowohl Hitze als auch Kälte, können zu Ermüdung, Muskelsteifheit oder anderen physiologischen Veränderungen führen, die sich auf die Bewegungsmuster auswirken.
	\item \textbf{Elektromagnetische Störungen}: Elektromagnetische Störungen können die Signalqualität der Sensoren beeinträchtigen. Diese Störungen können von verschiedenen Quellen wie elektronischen Geräten, Stromleitungen oder Funkwellen ausgehen. Um genaue und zuverlässige Daten zu gewährleisten, ist es wichtig, elektromagnetische Störungen zu minimieren.
\end{itemize}

Um die Genauigkeit und Zuverlässigkeit der Bewegungsanalyse zu gewährleisten und die oben genannten Umwelteinflüsse zu minimieren, ist es entscheidend, standardisierte Testbedingungen zu schaffen. Eine effektive Lösung besteht darin, die Tests in einem kontrollierten Umfeld durchzuführen, wie bereits bei der Vermeidung von Verzerrungen erwähnt. Beispielsweise kann ein Laufband in einem klimatisierten Raum verwendet werden, um konstante Bedingungen zu gewährleisten. Der Raum sollte gleichmäßig gekühlt sein, um Temperaturschwankungen zu vermeiden. Elektromagnetische Störungen sollten durch Abschirmungen minimiert werden. Um den Einfluss des Bewegungsziels zu kontrollieren, sollten den Testpersonen klare und einheitliche Anweisungen gegeben werden. Im Testraum sollten Markierungen platziert werden, um sicherzustellen, dass alle Testpersonen dieselben Bewegungsziele verfolgen. Durch diese Maßnahmen können konsistente und vergleichbare Daten gesammelt werden, die für die Entwicklung und Optimierung von Bewegungserkennungssystemen von großer Bedeutung sind.

\section{Hardware}
Eine der Hauptschwierigkeiten bei der Erfassung von Bewegungsdaten mit dem Arduino besteht in der Platzierung und dem Gewicht der Energieversorgung. Da der Arduino für mobile Anwendungen eine externe Stromquelle wie eine Powerbank benötigt, kann das zusätzliche Gewicht die natürlichen Bewegungsabläufe beeinträchtigen. Besonders beim Rennen kann die Powerbank störend wirken und die Bewegungsfreiheit einschränken. Ein erhöhtes Gewicht kann die Bewegungsfähigkeit der versuchenden erheblich beeinträchtigen \cite{monfrini2023technological} und somit zu verfälschten Messdaten führen, da die Bewegungen nicht mehr natürlich und ungehindert ausgeführt werden. Es ist daher wichtig, eine leichte und ergonomische Lösung für die Energieversorgung zu finden, um die Genauigkeit der Bewegungserfassung zu gewährleisten.

\section{Messfrequenz und Datenerfassung}
Eine zentrale Herausforderung ist die Festlegung der optimalen Messfrequenz. Eine zu niedrige Frequenz kann dazu führen, dass wichtige Bewegungsmuster nicht erfasst werden, während eine zu hohe Frequenz eine große Menge an Daten erzeugt, die schwer zu verarbeiten sind und die Systemleistung beeinträchtigen kann. Zusätzlich kann es durch Verzögerungen oder Signalstörungen zu Datenlücken kommen, die die Analyse verfälschen.